\documentclass[10pt,a4paper]{amsart}
\usepackage[margin=19mm]{geometry}
\usepackage{amsmath,amssymb,mathtools}
\usepackage[scr=rsfs,cal=euler]{mathalfa}
\usepackage{xfrac}
\usepackage[unicode,hidelinks,bookmarks=false]{hyperref}
\newcommand{\ie}{i.e.\ }
\newcommand{\cf}{cf.\ }
\newcommand{\resp}{respectively\ }
\newcommand{\Subsection}{\subsection}
\providecommand{\nobreakdash}{\nobreak\hbox{-}}
\setlength{\emergencystretch}{2em}
\multlinegap0pt
\newtheorem{theorem}{Theorem}[section]
\newtheorem{lemma}[theorem]{Lemma}
\newtheorem{conjecture}[theorem]{Conjecture}
\newtheorem{proposition}[theorem]{Proposition}
\newtheorem{corollary}[theorem]{Corollary}
\newtheorem{definition}[theorem]{Definition}
\newtheorem{example}[theorem]{Example}
\newtheorem{remark}[theorem]{Remark}
\begin{document}
\title[Maximum Forest Number of General Bipartite Graphs: Structura]{Maximum Forest Number of General Bipartite Graphs: Structural and Complexity Results}
\author{Daniel Iľkovič}
\date{}
\maketitle
\noindent\emph{Source and licence.} Maximum Forest Number of General Bipartite Graphs: Structural and Complexity Results. Version arXiv:2606.26761v1 (CC BY 4.0). This material is distributed under the Creative Commons Attribution 4.0 International licence, http://creativecommons.org/licenses/by/4.0/, as recorded in the arXiv OAI licence metadata for this exact version and shown on the abstract page https://arxiv.org/abs/2606.26761v1. Full rights notice: licenses/arxiv-2606.26761-CC-BY-4.0.txt.\par\smallskip
\noindent\textit{Editorial scope.} Counted unit: the original \texttt{theorem} environment \texttt{thm:main} at source lines 79--98 together with its complete proof at lines 100--183; the delivered document keeps source lines 75--184 so that every referenced label and every citation resolves inside it. Native editable source is the paper's own concatenated TeX; the AMS wrapper is editorial formatting only. No publisher class, upstream script or unrelated artwork is redistributed.
\section{Main results}

Our main theorem provides the exact maximum forest number over the family of valid bipartite graphs for any $m \le n$ and $d \ge 2$.


\begin{theorem} \label{thm:main}
Let $\mathcal{B}_{m,n,d}$ be the family of simple bipartite graphs $B = (X, Y)$ with $2 \le |X| = m \le n = |Y|$ that satisfy the Ore-type condition $\sigma_2(B) \ge d$, where $2 \le d \le 2m$.
The maximum forest number over this family, $\max_{B \in \mathcal{B}_{m,n,d}} f(B)$, is bounded by the following rules:

1. If $d = 2$, or if $d = 3$ and $m = n$, $\max f(B) = m + n$.

2. If $d \ge 4$, or if $d = 3$ and $m < n$, let $p = \lceil d/2 \rceil$ and $\delta = d \bmod 2$. Define $\alpha = m - p$, $\beta = n - p$, $\alpha' = \alpha + \delta$, and $\beta' = \beta + \delta$.
If $\beta = 0$, then $\max f(B) = n + 1$.
Otherwise ($\beta > 0$), calculate:
\[ c = \beta' - \alpha \beta - \alpha' \]
\[ l_0 = 1 + \frac{-c + \sqrt{c^2 + 4 \alpha \beta \beta'}}{2\beta} \]
\[ l_n = 1 + \left\lfloor \frac{\beta'}{p-1} \right\rfloor \]
Define the discrete candidate set of $X$-partition sizes:
\[ \mathcal{L} = \Big\{ 2, m, \alpha+2, l_n, \lfloor l_0 \rfloor, \lceil l_0 \rceil \Big\} \cap [2, m] \]
For each $l \in \mathcal{L}$, let $R(l) = \min \big( n, r_1(l), r_2(l) \big)$, where:
\[ r_1(l) = \beta + 1 + \left\lfloor \frac{\beta'}{l-1} \right\rfloor \]
\[ r_2(l) = \begin{cases} 1 + \left\lfloor \frac{\alpha'}{l - \alpha - 1} \right\rfloor & \text{if } l \ge \alpha + 2 \\ \infty & \text{if } l \le \alpha + 1 \end{cases} \]
Then:
\[ \max_{B \in \mathcal{B}_{m,n,d}} f(B) = \max \left( n+1, \max_{l \in \mathcal{L}} \big(l + R(l)\big) \right) \]
\end{theorem}

\begin{proof}
For $d=2$, we can construct a graph consisting of a disjoint union of $m-1$ independent edges and one star $K_{1, n-m+1}$.
This graph covers all $m+n$ vertices, is a forest, and has minimum degree 1, satisfying $\sigma_2(B) \ge 2$.
For $d=3$ and $m=n$, the path $P_{2n}$ covers all vertices, is a forest, and has degrees 1 and 2, yielding $\sigma_2(B) \ge 3$.
In both cases, $\max f(B) = m+n$.

For $d \ge 4$, or $d=3$ and $m < n$: The complete bipartite graph $K_{m,n}$ belongs to $\mathcal{B}_{m,n,d}$ (since the minimum degree is $m \ge \lceil d/2 \rceil$). Any induced forest in $K_{m,n}$ can contain at most one vertex from one of the partitions if it contains two or more from the other (to avoid a cycle of length 4). Thus, the maximum induced forest consists of all $n$ vertices from $Y$ and exactly $1$ vertex from $X$, yielding $f(K_{m,n}) = n+1$. Therefore, $\max f(B) \ge n+1$.

Suppose there exists a graph $B \in \mathcal{B}_{m,n,d}$ with an induced forest $S \subseteq V(B)$ of maximum size $s \ge n+2$.
Let $L = S \cap X$ and $R = S \cap Y$, with $|L| = l$ and $|R| = r$, forcing $l+r = s$.
To achieve $s \ge n+2$, we must have $l \ge 2$ and $r \ge 2$.

The condition $\sigma_2(B) \ge d$ guarantees that the sum of degrees inside $L$ satisfies $\sum_{x \in L} d(x) \ge lp - \delta$, where $p = \lceil d/2 \rceil$ and $\delta = d \bmod 2$.
To see this, consider the vertex of minimum degree in $L$, say $v_{min}$. Because $l \ge 2$, we can pair $v_{min}$ with another vertex $v_1 \in L$. The Ore-type condition requires $d(v_{min}) + d(v_1) \ge d = 2p - \delta$. Furthermore, at most one vertex in $L$ can have degree strictly less than $p$; if two distinct vertices had degrees $\le p-1$, their sum would be $\le 2p-2 < d$, a contradiction. Thus, the remaining $l-2$ vertices in $L$ must each have degree at least $p$. Summing these contributions, the total degree sum in $L$ is at least $(2p - \delta) + (l-2)p = lp - \delta$.

Each vertex in $L$ sends at most $n-r$ edges to $Y \setminus R$.
Counting edges strictly inside $B[S]$ from the perspective of $X$:
\[ e(B[S]) \ge \sum_{x \in L} d(x) - l(n-r) \ge lp - \delta - l(n-r) \]
Because $B[S]$ is a forest, it contains at most $s-1 = l+r-1$ edges.
Thus, we have the inequality:
\[ lp - \delta - l(n-r) \le l+r-1 \]
Substituting $\beta = n-p$ and rearranging to isolate $r$ yields:
\[ r \le \beta + 1 + \frac{\beta + \delta}{l-1} \]
Because $r$ is an integer, we can take the floor of the right side. Recalling that $\beta' = \beta + \delta$, we establish the first bound:
\[ r \le \beta + 1 + \left\lfloor \frac{\beta'}{l-1} \right\rfloor = r_1(l) \]

By symmetric logic regarding the $Y$-side degrees, we count edges strictly inside $B[S]$ from the perspective of $R$.
The sum of degrees in $R$ satisfies $\sum_{y \in R} d(y) \ge rp - \delta$.
Each vertex in $R$ sends at most $m-l$ edges to $X \setminus L$.
Thus, counting edges strictly inside $B[S]$ yields:
\[ e(B[S]) \ge rp - \delta - r(m-l) \]
Since $e(B[S]) \le l+r-1$, we have:
\[ rp - \delta - r(m-l) \le l+r-1 \]
Substituting $\alpha = m-p$ and rearranging to isolate $r$ yields, for $l \ge \alpha + 2$:
\[ r \le 1 + \frac{\alpha + \delta}{l - \alpha - 1} \]
Using the fact that $\alpha' = \alpha + \delta$ and $r$ must be an integer, we obtain the second bound:
\[ r \le 1 + \left\lfloor \frac{\alpha'}{l - \alpha - 1} \right\rfloor = r_2(l) \]
If $l \le \alpha + 1$, the term $l - \alpha - 1$ is non-positive, meaning the inequality places no upper bound on $r$, which corresponds to $r_2(l) = \infty$.
Therefore, any valid forest partitions must satisfy $r \le \min(n, r_1(l), r_2(l))$.

Before proceeding with the continuous optimization, consider the boundary case where $\beta = 0$. Since $\beta' = \delta$, the first bound evaluates to $r \le 1 + \lfloor \delta / (l-1) \rfloor$. Since $\delta \in \{0, 1\}$ and $l \ge 2$, this forces $r \le 2$. If $r \le 1$, then $s = l + r \le m + 1 \le n + 1$. If $r = 2$, this requires $\delta = 1$ and $l = 2$, giving $s = 4$. Since $d \ge 4$ or $m < n$ (so $n \ge 3$), we have $n+1 \ge 4$, thus $s \le n+1$. In either case, if $\beta = 0$, the maximum size is $n+1$, and we may assume $\beta > 0$ for the remainder of the analysis.

The objective is to maximize $l + \min(n, r_1(l), r_2(l))$ over $l \in [2, m]$.
The functions $r_1(l)$ and $r_2(l)$ are discrete versions of continuous hyperbolic curves, which are convex functions for $l > 1$. The sum $l + \min(n, r_1(l), r_2(l))$ is a continuous piecewise convex function on the interval $[2, m]$. The global maximum of a continuous piecewise convex function over a bounded interval must occur at the boundaries of its convex sub-intervals. These sub-intervals are delineated precisely by the global domain boundaries ($2$ or $m$) and the curve intersection points where the minimum function transitions. Because $r_2(l)$ has a discrete jump from $\infty$ to a finite value exactly at $l = \alpha + 2$, this point represents a hard boundary in the piecewise function and must be explicitly evaluated to capture any maximum that occurs at the asymptote jump.

The intersection of the continuous curve for $r_1(l)$ with the maximum possible value $n$ occurs when:
\[ \beta + 1 + \frac{\beta'}{l-1} = n \]
Substituting $\beta = n - p$ and solving for $l$ yields $l = 1 + \frac{\beta'}{p-1}$.
Taking the integer floor gives the discrete candidate $l_n = 1 + \left\lfloor \frac{\beta'}{p-1} \right\rfloor$.
The maximum valid value for the curve $r_2(l)$ occurs at $l = \alpha + 2$, where $r_2(\alpha+2) = 1 + \alpha' \le m \le n - 1 < n$. Because the curve $r_2(l)$ is bounded strictly below $n$ on its valid domain, the intersection $r_2(l) = n$ does not exist.

The intersection of the continuous curves for $r_1(l)$ and $r_2(l)$ occurs when:
\[ \beta + 1 + \frac{\beta'}{l-1} = 1 + \frac{\alpha'}{l - \alpha - 1} \]
Substituting $u = l-1$ and clearing denominators yields the quadratic equation:
\[ \beta u^2 + (\beta' - \alpha \beta - \alpha')u - \alpha \beta' = 0 \]
Recall our definition $c = \beta' - \alpha \beta - \alpha'$. Applying the quadratic formula to find the positive root for $u$ and substituting back $l = u + 1$ yields the exact continuous intersection point:
\[ l_0 = 1 + \frac{-c + \sqrt{c^2 + 4 \alpha \beta \beta'}}{2\beta} \]
Since $l$ must be an integer, the maximum over the discrete domain is bounded by evaluating the function at the boundary points $2$ and $m$, and the nearest integers to the continuous intersection points $l_n$ and $l_0$. Thus, evaluating the set $\mathcal{L}$ covers all possible discrete extrema.

To prove sharpness, let $l \in \mathcal{L}$ and $r = R(l)$ be the optimizing pair yielding the maximum $s = l+r$.
We construct a valid bipartite graph $B = (X, Y)$ that attains this bound.
Let $S = L \cup R$ be the target forest with $L = S \cap X$ and $R = S \cap Y$, such that $|L| = l$ and $|R| = r$. 
Let the remaining vertices be $X_0 = X \setminus L$ and $Y_0 = Y \setminus R$.

To construct the acyclic graph $B[S]$, explicitly take a path $P$ that alternates between vertices in $L$ and $R$. 
If $l \ne r$, for any remaining unconnected vertices in the larger of the two sets ($L$ or $R$), connect them as leaves to the path such that the degrees inside $B[S]$ are distributed as evenly as possible. 
This explicit construction ensures that $B[S]$ is a simple path or a star forest (which is acyclic) containing exactly $l+r-1$ edges. 
Consequently, the average degree in $B[S]$ for vertices in $L$ is $\frac{l+r-1}{l} = 1 + \frac{r-1}{l}$, and for vertices in $R$ is $\frac{l+r-1}{r} = 1 + \frac{l-1}{r}$. 

Then, to ensure all vertices meet the minimum degree bounds ($p$ or $d-p$) required by the Ore-type condition $\sigma_2(B) \ge d$, explicitly add all possible edges between $L$ and $Y_0$, all possible edges between $R$ and $X_0$, and all possible edges between $X_0$ and $Y_0$. 
Because we only care about the lower bounds for the degrees to satisfy the theorem, adding all these cross edges is sufficient and does not break the acyclic property of $B[S]$ (since there are no edges among $X_0$ and $Y_0$ that could form a cycle back into $S$).

Let us verify the degrees of the vertices in this construction. 
Every vertex $u \in X_0$ connects to all $r$ vertices in $R$ and all $n-r$ vertices in $Y_0$, achieving degree $n$. Since $n \ge \lceil d/2 \rceil = p$, we have $d(u) \ge p$.
Every vertex $v \in Y_0$ connects to all $l$ vertices in $L$ and all $m-l$ vertices in $X_0$, achieving degree $m$. Since $m \ge p$, we have $d(v) \ge p$.
For a vertex $x \in L$, its degree is exactly its degree in $B[S]$ plus its $n-r$ connections to $Y_0$. 
Recall the bound $r \le r_1(l) = \beta + 1 + \lfloor \frac{\beta'}{l-1} \rfloor$. As shown earlier, this algebraic inequality is equivalent to $lp - \delta - l(n-r) \le l+r-1$, which implies the average degree inside $B[S]$ for $L$ is at least $p - (n-r) - \frac{\delta}{l}$. By evenly distributing the edges in $B[S]$, the minimum degree in $L$ is at least $p - (n-r)$ (with at most one vertex potentially having degree $p - (n-r) - 1$ if $\delta=1$). Thus, after adding the $n-r$ edges to $Y_0$, every $x \in L$ has $d(x) \ge p$ (or $d-p$ for at most one vertex). 
By a symmetric argument, the bound $r \le r_2(l)$ guarantees that for any $y \in R$, the average degree inside $B[S]$ is at least $p - (m-l) - \frac{\delta}{r}$. Adding the $m-l$ edges to $X_0$ ensures $d(y) \ge p$ (or $d-p$ for at most one vertex).

Consequently, any pair of vertices $x_1, x_2 \in X$ will have a degree sum $d(x_1) + d(x_2) \ge p + (d-p) = d$. Similarly, any pair in $Y$ will have a degree sum of at least $d$. 
Thus, the Ore-type condition $\sigma_2(B) \ge d$ is satisfied. 
Meanwhile, the induced subgraph $B[S]$ contains exactly the edges of the initial explicit construction and no more, making it a valid induced forest of size $s = l+r$.
This proves the exactness of the bound.
\end{proof}


\end{document}
